\documentclass[11pt]{article}
% Shared preamble for the pyMack LST documentation set.
% Used by: mean_boundary_layer.tex, disturbance_equations.tex,
%          numerical_methods.tex, validation.tex
% Each document begins with:  \documentclass[11pt]{article}% Shared preamble for the pyMack LST documentation set.
% Used by: mean_boundary_layer.tex, disturbance_equations.tex,
%          numerical_methods.tex, validation.tex
% Each document begins with:  \documentclass[11pt]{article}% Shared preamble for the pyMack LST documentation set.
% Used by: mean_boundary_layer.tex, disturbance_equations.tex,
%          numerical_methods.tex, validation.tex
% Each document begins with:  \documentclass[11pt]{article}\input{lst_preamble}
\usepackage[margin=1in]{geometry}
\usepackage{amsmath,amssymb,mathtools}
\usepackage{booktabs}
\usepackage{enumitem}
\usepackage{array}
\usepackage{graphicx}
\usepackage{caption}
\usepackage{xcolor}
\usepackage[hidelinks]{hyperref}

\graphicspath{{figures/}}
\captionsetup{font=small,labelfont=bf}

\definecolor{kicol}{RGB}{31,143,107}
\definecolor{boxbg}{RGB}{238,246,242}
\definecolor{rmkbg}{RGB}{244,244,244}

\newcommand{\plain}[1]{\par\smallskip\noindent{\color{kicol}\textbf{In plain words.}}\ \textit{#1}\par\smallskip}
\newcommand{\problembox}[1]{%
  \par\smallskip\noindent\fcolorbox{kicol}{boxbg}{%
  \parbox{0.965\linewidth}{\smallskip\textbf{\color{kicol}What problem is solved here.}\ \ #1\smallskip}}\par\smallskip}
\newcommand{\remark}[2]{%
  \par\smallskip\noindent\colorbox{rmkbg}{%
  \parbox{0.965\linewidth}{\smallskip\footnotesize\textbf{#1}\ \ #2\smallskip}}\par\smallskip}
\newcommand{\bridge}[1]{\par\noindent\textit{#1}\par\medskip}

\newcommand{\D}{\mathrm{D}}
\newcommand{\imag}{\mathrm{i}}
\newcommand{\Real}{\operatorname{Re}}
\newcommand{\Imag}{\operatorname{Im}}
\newcommand{\Mae}{M_e}
\newcommand{\Ree}{\mathit{Re}}
\renewcommand{\bar}[1]{\overline{#1}}
\newcommand{\hatu}{\hat{u}}
\newcommand{\hatv}{\hat{v}}
\newcommand{\hatT}{\hat{T}}
\newcommand{\hatp}{\hat{p}}
\newcommand{\hatq}{\hat{q}}
\newcommand{\nuR}{\nu_{\!R}}

\usepackage[margin=1in]{geometry}
\usepackage{amsmath,amssymb,mathtools}
\usepackage{booktabs}
\usepackage{enumitem}
\usepackage{array}
\usepackage{graphicx}
\usepackage{caption}
\usepackage{xcolor}
\usepackage[hidelinks]{hyperref}

\graphicspath{{figures/}}
\captionsetup{font=small,labelfont=bf}

\definecolor{kicol}{RGB}{31,143,107}
\definecolor{boxbg}{RGB}{238,246,242}
\definecolor{rmkbg}{RGB}{244,244,244}

\newcommand{\plain}[1]{\par\smallskip\noindent{\color{kicol}\textbf{In plain words.}}\ \textit{#1}\par\smallskip}
\newcommand{\problembox}[1]{%
  \par\smallskip\noindent\fcolorbox{kicol}{boxbg}{%
  \parbox{0.965\linewidth}{\smallskip\textbf{\color{kicol}What problem is solved here.}\ \ #1\smallskip}}\par\smallskip}
\newcommand{\remark}[2]{%
  \par\smallskip\noindent\colorbox{rmkbg}{%
  \parbox{0.965\linewidth}{\smallskip\footnotesize\textbf{#1}\ \ #2\smallskip}}\par\smallskip}
\newcommand{\bridge}[1]{\par\noindent\textit{#1}\par\medskip}

\newcommand{\D}{\mathrm{D}}
\newcommand{\imag}{\mathrm{i}}
\newcommand{\Real}{\operatorname{Re}}
\newcommand{\Imag}{\operatorname{Im}}
\newcommand{\Mae}{M_e}
\newcommand{\Ree}{\mathit{Re}}
\renewcommand{\bar}[1]{\overline{#1}}
\newcommand{\hatu}{\hat{u}}
\newcommand{\hatv}{\hat{v}}
\newcommand{\hatT}{\hat{T}}
\newcommand{\hatp}{\hat{p}}
\newcommand{\hatq}{\hat{q}}
\newcommand{\nuR}{\nu_{\!R}}

\usepackage[margin=1in]{geometry}
\usepackage{amsmath,amssymb,mathtools}
\usepackage{booktabs}
\usepackage{enumitem}
\usepackage{array}
\usepackage{graphicx}
\usepackage{caption}
\usepackage{xcolor}
\usepackage[hidelinks]{hyperref}

\graphicspath{{figures/}}
\captionsetup{font=small,labelfont=bf}

\definecolor{kicol}{RGB}{31,143,107}
\definecolor{boxbg}{RGB}{238,246,242}
\definecolor{rmkbg}{RGB}{244,244,244}

\newcommand{\plain}[1]{\par\smallskip\noindent{\color{kicol}\textbf{In plain words.}}\ \textit{#1}\par\smallskip}
\newcommand{\problembox}[1]{%
  \par\smallskip\noindent\fcolorbox{kicol}{boxbg}{%
  \parbox{0.965\linewidth}{\smallskip\textbf{\color{kicol}What problem is solved here.}\ \ #1\smallskip}}\par\smallskip}
\newcommand{\remark}[2]{%
  \par\smallskip\noindent\colorbox{rmkbg}{%
  \parbox{0.965\linewidth}{\smallskip\footnotesize\textbf{#1}\ \ #2\smallskip}}\par\smallskip}
\newcommand{\bridge}[1]{\par\noindent\textit{#1}\par\medskip}

\newcommand{\D}{\mathrm{D}}
\newcommand{\imag}{\mathrm{i}}
\newcommand{\Real}{\operatorname{Re}}
\newcommand{\Imag}{\operatorname{Im}}
\newcommand{\Mae}{M_e}
\newcommand{\Ree}{\mathit{Re}}
\renewcommand{\bar}[1]{\overline{#1}}
\newcommand{\hatu}{\hat{u}}
\newcommand{\hatv}{\hat{v}}
\newcommand{\hatT}{\hat{T}}
\newcommand{\hatp}{\hat{p}}
\newcommand{\hatq}{\hat{q}}
\newcommand{\nuR}{\nu_{\!R}}

\usepackage{tikz}
\usetikzlibrary{arrows.meta,positioning}

\title{\textbf{The Big-Picture Map: How the Stages of the \texttt{pyMack} Workflow Fit Together}\\[3pt]
\large What each stage consumes, what it produces, which formulation and solver is used when,\\
and how the chain ends at a transition estimate}
\author{\texttt{pyMack} technical documentation}
\date{\today}

\begin{document}
\maketitle
\tableofcontents
\bigskip

%=====================================================================
\section{Scope}\label{sec:scope}
%=====================================================================
The five companion documents each descend into one stage of the linear stability workflow: the
steady base flow (\emph{Mean Boundary Layer}), the linearised operator and its discretisation
(\emph{Derivation of the Disturbance Equations}), the spectral eigenvalue algorithms
(\emph{Numerical Methods}), the independent local solver (\emph{Shooting Solver}), and the
evidence (\emph{Validation}). What none of them states in one place is the connective tissue: at which
stage the boundary-layer profile is actually consumed, which of the two stability formulations
(temporal or spatial) is used for which product, when the global eigenvalue solver is preferred
over the local shooting solver, and how the chain of intermediate objects terminates in a
neutral curve and an $N$-factor. This document supplies that map. Section~\ref{sec:pipeline}
presents the pipeline end to end; Section~\ref{sec:baseflow} traces every consumer of the mean
flow; Section~\ref{sec:formulations} states the division of labour between the temporal and
spatial formulations; Section~\ref{sec:engines} does the same for the four eigenvalue engines;
Section~\ref{sec:deliverables} follows the eigenvalues into neutral curves, $N$-factors, and
dimensional predictions; Section~\ref{sec:mach6} walks the canonical Mach~6 production case
through every stage; Section~\ref{sec:validationmap} locates the validation layers on the map.
Nothing here is new material: every statement is a cross-reference into the companion documents
and the code, cited by file.

%=====================================================================
\section{The pipeline at a glance}\label{sec:pipeline}
%=====================================================================
Everything the framework produces is assembled from one repeated primitive: an eigenvalue solve
of the linearised disturbance equations at a single parameter point. The pipeline is therefore a
chain of \emph{preparations} for that solve (stages 1--2), the \emph{solve itself} in one of two
formulations by one of several engines (stage 3), an \emph{acceptance test} that separates the
physical mode from discretisation artifacts (stage 4), and \emph{aggregation} of many accepted
solves into the engineering products (stages 5--6). For a normal-mode disturbance
\cite{mack1984,schmid2001}
\begin{equation}\label{eq:ansatz}
  q'(x,y,z,t)=\hatq(y)\,e^{\imag(\alpha x+\beta z-\omega t)},
\end{equation}
the two formulations differ only in which of $(\alpha,\omega)$ is prescribed real and which is
the complex unknown:
\begin{equation}\label{eq:formulations}
  \text{temporal: given }\alpha\in\mathbb R,\ \text{find }c=\omega/\alpha\in\mathbb C
  \ (\omega_i=\alpha c_i);\qquad
  \text{spatial: given }\omega\in\mathbb R,\ \text{find }\alpha\in\mathbb C\ (\sigma=-\alpha_i).
\end{equation}

\begin{figure}[H]\centering
\begin{tikzpicture}[
  node distance=5.5mm and 9mm,
  stage/.style={draw, rounded corners=2pt, align=center, fill=boxbg, inner sep=5pt,
                font=\small, minimum width=34mm},
  side/.style={draw, rounded corners=2pt, align=center, fill=rmkbg, inner sep=4pt,
               font=\footnotesize},
  lbl/.style={font=\footnotesize\itshape, text=black!60},
  arr/.style={-{Stealth[length=2.6mm]}, thick}
]
\node[stage] (cond) {1. Flow conditions\\ \footnotesize $\Mae$, wall BC, gas model, $T_e$};
\node[stage, below=of cond] (bl) {2. Mean boundary layer\\ \footnotesize self-similar BVP:
  $\bar U,\bar T,\bar\rho,\bar\mu,\bar\kappa$ + derivatives\\ \footnotesize \texttt{baseflow.py}, \texttt{boundary\_layer.py}};
\node[stage, below=of bl] (ops) {3. Disturbance operator at one $(R,\alpha\ \text{or}\ \omega)$\\
  \footnotesize temporal GEVP $A\boldsymbol\phi=cB\boldsymbol\phi$ \ / \ spatial QEP
  $(C_0{+}\alpha C_1{+}\alpha^2C_2)\boldsymbol\phi=0$\\
  \footnotesize \texttt{equations.py}, \texttt{temporal\_solver.py}, \texttt{solver.py}};
\node[stage, below=of ops] (accept) {4. Mode acceptance\\ \footnotesize freestream decay +
  domain stationarity};
\node[stage, below=of accept] (sweep) {5. Sweeps\\ \footnotesize growth curves
  $\omega_i(\alpha)$, $\sigma(\omega)$;\ neutral curves ($\omega_i{=}0$ or $\sigma{=}0$)\\
  \footnotesize \texttt{analysis.py}};
\node[stage, below=of sweep] (nfac) {6. $N$-factor and transition\\ \footnotesize
  $N(R)=\int\max(\sigma_L,0)\,\mathrm dR$;\ $e^N$; $N\approx9$--$11$};
\node[side, right=of ops, yshift=-1mm] (shoot) {Shooting solver\\ (refine / verify one mode)\\
  \texttt{temporal\_shooting.py}, \texttt{mack\_shooting.py}};
\node[side, right=of nfac] (dim) {Dimensionalisation\\ kHz, mm, 1/m\\ \texttt{scales.py};
  cone: \texttt{cone.py}};
\draw[arr] (cond) -- (bl);
\draw[arr] (bl) -- node[lbl, right=1mm] {profile sampled into every coefficient matrix} (ops);
\draw[arr] (ops) -- node[lbl, right=1mm] {full spectrum or targeted mode} (accept);
\draw[arr] (accept) -- node[lbl, right=1mm] {one trusted eigenvalue per parameter point} (sweep);
\draw[arr] (sweep) -- node[lbl, right=1mm] {spatial $\sigma_L$ at fixed $F$} (nfac);
\draw[arr, <->] (ops) -- node[lbl, above=0.5mm, align=center] {} (shoot);
\draw[arr] (nfac) -- (dim);
\end{tikzpicture}
\caption{The \texttt{pyMack} pipeline. The mean flow (stage 2) is computed once per flow
condition and consumed by every operator assembly (stage 3); stages 3--4 constitute the repeated
eigenvalue primitive; stages 5--6 aggregate thousands of such solves into the engineering
products. The shooting solver sits beside stage 3 as an independent second route to the same
eigenvalue.}\label{fig:pipeline}
\end{figure}

Table~\ref{tab:stages} states, for each stage, what enters, what leaves, and who consumes the
output. The remainder of the document expands each row.

\begin{table}[H]\centering\small
\renewcommand{\arraystretch}{1.3}
\begin{tabular}{@{}>{\raggedright\arraybackslash}p{2.6cm}>{\raggedright\arraybackslash}p{3.4cm}
>{\raggedright\arraybackslash}p{4.2cm}>{\raggedright\arraybackslash}p{4.2cm}@{}}
\toprule
\textbf{Stage} & \textbf{Input} & \textbf{Output} & \textbf{Consumed by} \\
\midrule
Mean flow & $\Mae$, wall BC, transport law & profiles $\bar U,\bar T,\bar\rho,\bar\mu,\bar\kappa$
  and $y$-derivatives on a similarity grid; $\delta^*/L^*$ & every coefficient matrix of stage 3;
  scale conversions; $y_{\max}$ sizing \\
Operator assembly & profiles sampled on the collocation grid; $(R,\alpha)$ or $(R,\omega)$ &
  matrix pencil $(A,B)$ or QEP blocks $(C_0,C_1,C_2)$; pointwise $6\times6$ $A(y)$ for shooting &
  eigenvalue engines \\
Eigenvalue solve & the operator + (spatial only) a shift target or seed & raw spectrum, or one
  refined eigenvalue & mode acceptance \\
Mode acceptance & raw eigenvalues + eigenfunctions & the discrete boundary-layer mode &
  growth curves, neutral tracing \\
Sweeps & accepted modes over $(R,\alpha)$ or $(R,\omega)$ grids & $\omega_i$ or $\sigma_L$
  fields; neutral branches & $N$-factor; overlays; verification verdicts \\
$N$-factor & $\sigma_L(R)$ at fixed $F$ & $N(R)$, $A/A_0=e^N$ & transition estimate;
  dimensional plots \\
\bottomrule
\end{tabular}
\caption{Stage-by-stage data flow. One mean-flow solve serves an entire stability survey;
one eigenvalue solve serves exactly one parameter point.}\label{tab:stages}
\end{table}

%=====================================================================
\section{Stage 1--2: the mean boundary layer and its consumers}\label{sec:baseflow}
%=====================================================================
The base flow answers a question that must be settled before any stability problem exists: about
\emph{what} state are the equations linearised? The self-similar compressible flat-plate solution
(\emph{Mean Boundary Layer}, Sections 2--3; \texttt{pymack/baseflow.py}) provides
$\bar U(y),\bar T(y),\bar\rho=1/\bar T$, the transport properties
$\bar\mu(\bar T),\bar\kappa(\bar T)$ with their temperature derivatives, and the first and second
$y$-derivatives of everything, on a wall-normal grid.

\paragraph{The profile is computed once and reused everywhere.} Because the solution is
self-similar, the nondimensional profile depends only on $(\Mae,\ \text{wall condition},\
\text{transport law},\ T_e)$ and \emph{not} on the streamwise station: the Reynolds number $R$
enters the disturbance operator as an explicit parameter (through the viscous prefactor
$\nuR=1/(\Ree\bar\rho)$), not through the profile. One call to
\texttt{flat\_plate}/\texttt{make\_mack\_profile}/\texttt{generate\_boundary\_layer} therefore
serves an entire $(R,\alpha)$ or $(R,\omega)$ survey of thousands of eigenvalue solves. This is
the single most economical fact of the workflow: stage 2 costs one boundary-value solve
\cite{kierzenka2001}, stage 3--4 cost one dense eigensolve \emph{per parameter point}.

\paragraph{Consumer 1: the coefficient matrices.} The profile's role in the eigenvalue problem
is to populate coefficients. Every entry of the temporal blocks $A,B$
(\texttt{temporal\_solver.py}), of the spatial blocks $C_0,C_1,C_2$ (\texttt{equations.py}), and
of the pointwise shooting matrix $A(y)$ (\texttt{mack\_shooting.py}) is an algebraic combination
of $\bar U,\bar U',\bar U''$, $\bar T,\bar T',\bar T''$, $\bar\rho$, $\bar\mu$,
$\mathrm d\mu/\mathrm dT$, $\mathrm d^2\mu/\mathrm dT^2$, $\bar\kappa$ and its derivatives,
evaluated at the grid points (\emph{Numerical Methods}, Section 2; \emph{Shooting Solver},
Section 2). The profile is
sampled onto the Chebyshev collocation grid (or the RK4 marching grid) by the single shared
resampler \texttt{scales.sample\_baseflow}; the solvers never re-solve the mean flow.

\paragraph{Consumer 2: the length scales.} The stability solvers store profiles on a
$y/\delta^*$ grid, while Mack's figures and all neutral-curve axes use the Falkner--Skan scale
$L^*=\sqrt{\nu_ex/U_e}$ with $R=\sqrt{Re_x}$ \cite{mack1984}. The conversion factor
$\delta^*/L^*$ is a property of the profile, computed by quadrature of $(\bar\rho\bar U)$ deficit
across the layer, and every $\alpha_{\delta^*}\!\leftrightarrow\!\alpha_L$ or
$y_{\delta^*}\!\leftrightarrow\!y_L$ conversion in the repository routes through
\texttt{pymack/scales.py} (\emph{Mean Boundary Layer}, Section 4).

\paragraph{Consumer 3: the domain height.} The collocation domain is truncated at
$y_{\max}=m\,(\delta^*/L^*)$ with $m\approx8$--$12$ for the wall-trapped second mode and
$m\approx35$--$45$ for the diffuse first mode (\emph{Numerical Methods}, Section 2), so even the
grid geometry of stage 3 is set by a boundary-layer thickness that stage 2 supplies.

\paragraph{What the profile is \emph{not} used for.} The $N$-factor integration and the
dimensionalisation consume only eigenvalue outputs ($\sigma_L$ vs $R$) and an edge state
($U_e,\nu_e,T_e$); the profile does not reappear downstream of stage 3. If two runs share a
profile and differ in $R$, their difference is carried entirely by the operator, not by the base
flow.

%=====================================================================
\section{Stage 3, first decision: temporal or spatial formulation}\label{sec:formulations}
%=====================================================================
The repository genuinely uses \emph{both} formulations of Eq.~\eqref{eq:formulations}, not as
redundant alternatives but with a clean division of labour that follows from what each returns
and what each costs.

\paragraph{Why the temporal problem is the survey tool.} With $\alpha$ real and $c$ the
eigenvalue, the discretised problem is the \emph{linear} generalised eigenvalue problem
\begin{equation}\label{eq:gevp}
  \boxed{\;A(\alpha)\,\boldsymbol\phi=c\,B(\alpha)\,\boldsymbol\phi\;}
\end{equation}
(\emph{Numerical Methods}, Section 3), solved by QZ \cite{molerstewart1973} for \emph{all} $4n$
eigenvalues at once with no initial guess. This no-guess property is the reason the temporal
formulation is the discovery and survey tool \cite{malik1990}: when nothing is yet known about
the spectrum at a new $(\Mae,R)$, a temporal solve finds the second mode (phase speed
$c_r\approx1-1/\Mae$), the first mode ($c_r\approx0.3$--$0.6$), and the continuous-spectrum
clusters simultaneously. Temporal sweeps produce the growth maps $\omega_i(R,\alpha)$, the
temporal neutral curves, and the validation comparisons against Mack's temporal tables (Table
10.1 growth rates) and \"Ozgen \& K{\i}rcal{\i}'s Fig.~3 lobes \cite{mack1984,ozgen2008}.

\paragraph{Why the spatial problem is the transition tool.} A physical disturbance source
(roughness, acoustic forcing) oscillates at a fixed real frequency and grows in \emph{space}, so
transition prediction needs $\sigma=-\alpha_i$ at prescribed real $\omega$; the $N$-factor
integral is a streamwise integral of spatial growth and cannot be built from temporal quantities
without an approximate group-velocity conversion \cite{gaster1962}. With $\alpha$ as the
eigenvalue the viscous terms make the problem \emph{quadratic},
\begin{equation}\label{eq:qep}
  \big(C_0+\alpha C_1+\alpha^2C_2\big)\boldsymbol\phi=0,
\end{equation}
linearised to companion form at twice the size and solved near a shift target
$\alpha_0=\omega/c_{\mathrm{ph}}$ \cite{tisseur2001} (\emph{Numerical Methods}, Section 4). The
spatial formulation therefore requires an idea of where the mode lies; it is used once the mode
family is known, to produce the quantities transition actually consumes: fixed-frequency growth
curves $\sigma_L(R)$, spatial neutral branches, and $N$-factors.

\paragraph{The bridge between the two.} The two formulations meet in
\texttt{solve\_spatial\_from\_temporal} (\texttt{solver.py}): a temporal solve at
$\alpha_r=\omega/c_{\mathrm{est}}$ supplies the mode and a Gaster-type seed
$\alpha_i\approx-\alpha_rc_i/c_r$ \cite{gaster1962}, and Newton iteration on the dispersion
residual $F(\alpha)=\alpha\,c(\alpha)-\omega=0$ converges to the spatial eigenvalue. The temporal
solver's no-guess robustness thus bootstraps the spatial solver's physically relevant answer;
both spectral spatial routes return the same $\sigma$ (\emph{Numerical Methods}, Section 4).

\begin{table}[H]\centering\small
\renewcommand{\arraystretch}{1.3}
\begin{tabular}{@{}>{\raggedright\arraybackslash}p{4.0cm}>{\raggedright\arraybackslash}p{5.0cm}
>{\raggedright\arraybackslash}p{5.6cm}@{}}
\toprule
 & \textbf{Temporal (real $\alpha$, complex $c$)} & \textbf{Spatial (real $\omega$, complex $\alpha$)} \\
\midrule
Algebraic form & linear GEVP, Eq.~\eqref{eq:gevp} & quadratic QEP, Eq.~\eqref{eq:qep} \\
Needs a guess? & no (full spectrum by QZ) & yes (shift target or temporal seed) \\
Growth rate & $\omega_i=\alpha c_i$ & $\sigma=-\alpha_i$ \\
Used in this repository for & mode discovery and surveys; growth maps; temporal neutral curves;
  Mack Table 10.1 and Fig.~10.6 checks; \"Ozgen Fig.~3 lobes & fixed-frequency growth curves;
  spatial neutral branches; $N$-factors; Malik Case 6 anchor; Mach 6 production case;
  Mach 5.35 collaborator benchmark \\
Typical entry points & \texttt{temporal\_mode}, \texttt{temporal\_growth\_curve},
  \texttt{trace\_temporal\_neutral\_curve} & \texttt{spatial\_mode},
  \texttt{spatial\_growth\_curve}, \texttt{trace\_spatial\_neutral\_curve} \\
\bottomrule
\end{tabular}
\caption{Division of labour between the two formulations. The temporal problem is the
no-guess survey instrument; the spatial problem produces what transition prediction consumes.}
\label{tab:formulations}
\end{table}

%=====================================================================
\section{Stage 3, second decision: which eigenvalue engine}\label{sec:engines}
%=====================================================================
Within a formulation, \texttt{pyMack} keeps several numerical routes to the same eigenvalue,
deliberately: independent discretisations that agree on a mode constitute a verification
statement no single method can make \cite{malik1990}. The four engines and their roles:

\paragraph{Engine 1: temporal spectral (global).} Chebyshev collocation
\cite{orszag1971,trefethen2000} plus dense QZ on Eq.~\eqref{eq:gevp}
(\texttt{temporal\_solver.py} in \"Ozgen's temperature-form energy arrangement,
\texttt{solver.py} in Mack's enthalpy form; the two arrangements bracket the closure choice and
agree on the physical mode to discretisation accuracy). Chosen whenever the spectrum is unknown:
surveys, first looks at a new condition, neutral-curve sweeps.

\paragraph{Engine 2: spatial spectral (global, targeted).} The companion linearisation of
Eq.~\eqref{eq:qep} with shift--invert selection near $\alpha_0$, or the Newton route from a
temporal seed (\texttt{solver.py}). Chosen for every product that needs $\sigma$ at fixed
frequency, which is all of stage 6.

\paragraph{Engine 3: shooting (local).} The classical initial-value route
\cite{mack1984,ozgen2008}: the disturbance equations recast as a first-order system
$\mathrm dY/\mathrm dy=A(y)Y$ in a six-component state, the three freestream-decaying solutions
of the constant-coefficient limit $A_\infty$ marched from $y_{\max}$ to the wall by RK4 with QR
re-orthonormalisation after every step \cite{conte1966}, and the eigenvalue located by driving
the smallest singular value of the $3\times3$ wall matrix to zero,
\begin{equation}\label{eq:sigmamin}
  \sigma_{\min}\big(M(c)\big)\to0\quad\Longleftrightarrow\quad c\ \text{is an eigenvalue},
\end{equation}
by Nelder--Mead search \cite{neldermead1965} (\emph{Shooting Solver} document;
\texttt{temporal\_shooting.py}, \texttt{mack\_shooting.py}). The shooting engine returns
\emph{one} eigenvalue and needs a good guess, so it is never the survey tool. It is used in
exactly three situations:
\begin{enumerate}[nosep,leftmargin=1.8em]
  \item \textbf{Independent verification.} Agreement between spectral and shooting eigenvalues
        at the same conditions is the repository's strongest internal check that a mode is
        physics rather than discretisation artifact
        (\texttt{validation/test\_temporal\_shooting\_cross\_check.py}).
  \item \textbf{Refinement.} A spectral eigenvalue accurate to discretisation error is polished
        to tight tolerance by the local search seeded with the spectral value.
  \item \textbf{Regimes where the reduced EVP misses a branch.} For the low-/mid-Mach oblique
        first mode, the exact first-order shooting system (Mack's Appendix A operator, extended
        to the sixth- and eighth-order oblique systems) resolves the Table 10.1 growth rates to
        0.07--0.91\%, while the reduced temporal EVP returns a different, unphysical branch
        family; the shooting path is authoritative there (\emph{Validation}; the Mack Table 10.1
        machinery in \texttt{mack\_table\_10\_1.py} and
        \texttt{analysis.py:temporal\_growth\_scan\_3d\_shooting}).
\end{enumerate}

\paragraph{Engine 4: the dense spatial backend.} \texttt{pymack/dense.py} is a compact,
independently written spatial QEP path with its own Lees--Dorodnitsyn base flow and
most-amplified-seed curve tracking. It exists as an \emph{independent reference implementation}
for the Mach 6 second-mode production work: the cross-method consistency gate
(\texttt{validation/test\_spatial\_cross\_method\_consistency.py}) holds the dense backend and
the main \texttt{solve\_spatial} family to $|\Delta\alpha_r|\le2\times10^{-4}$,
$|\Delta\sigma|\le10^{-4}$ at the canonical Mach 6 point. The production Mach 6 runner uses it
as its default backend (Section~\ref{sec:mach6}).

\paragraph{Stage 4, always: mode acceptance.} Whatever the engine, a raw spectral solve returns
$4n$ eigenvalues, most of them the finite-grid rendering of the continuous spectrum. Two tests
select the discrete boundary-layer mode (\emph{Numerical Methods}, Section 5): the eigenfunction
must decay toward the freestream (edge-to-peak amplitude below $0.1$), and the
eigenvalue must be stationary under a change of domain height. The facade applies the decay test
always and the stationarity test automatically when no user guess anchors the branch
(\texttt{api.py}). The worked example of \emph{Numerical Methods} Section 6 shows why the tests
are not optional: at $\Mae=6$, $R=5500$ the largest-growth raw eigenvalue is spurious, and the
physical Mack mode is the \emph{second} entry of the sorted list.

%=====================================================================
\section{Stages 5--6: from eigenvalues to deliverables}\label{sec:deliverables}
%=====================================================================
Every engineering product is an aggregation of accepted stage-4 eigenvalues.

\paragraph{Growth curves.} At fixed $R$: sweep $\alpha$ and record $\omega_i=\alpha c_i$
(temporal, \texttt{temporal\_growth\_curve}) or sweep $\omega$ and record $\sigma=-\alpha_i$
(spatial, \texttt{spatial\_growth\_curve}). These are the per-station instability portraits and
the unit of comparison for most literature checks.

\paragraph{Neutral curves.} The neutral curve is the zero-growth contour of a growth field over
a two-parameter sweep: $\omega_i(R,\alpha)=0$ or $\sigma(R,\omega)=0$. It is traced by sweeping
the grid, applying stage 4 at every node, and extracting the zero contour, with secant or Brent
refinement of the crossings and continuation for branches that approach the continuous spectrum
(\texttt{trace\_temporal\_neutral\_curve}, \texttt{trace\_spatial\_neutral\_curve},
\texttt{analysis.py}). Each fixed $R$ inside the unstable pocket yields two crossings, the lower
(branch I) and upper (branch II) neutral points; the critical Reynolds number is the nose where
the two branches meet (\texttt{critical\_reynolds\_by\_max\_growth}).

\paragraph{$N$-factor.} The transition estimate follows a \emph{fixed physical frequency}
$F=\omega_L/R$ downstream: the spatial problem is solved at successive stations $R$, and the
accumulated amplification is
\begin{equation}\label{eq:nfactor}
  N(R)=\int_{R_0}^{R}\max(\sigma_L,0)\,\mathrm dR',\qquad \frac{A(R)}{A(R_0)}=e^{N(R)},
\end{equation}
integrated by the trapezoid rule from the lower neutral point $R_0$ of that frequency
(\texttt{integrate\_n\_factor}, \texttt{analysis.py}). Transition is empirically correlated with
$N\approx9$--$11$ in flight and $N\approx4$--$6$ in noisy conventional tunnels, the $e^N$ method
of Smith \& Gamberoni and van Ingen \cite{smith1956,vaningen1956}. The envelope over frequencies
of $N(R)$ curves is the transition-prediction product. Note the chain of dependencies: the
$N$-factor needs spatial growth rates at fixed $F$ (stage 5, spatial formulation, engines 2 or
4), which need accepted eigenvalues (stage 4), which need the operator (stage 3), which needs
the profile (stage 2). A single wrong stage poisons everything downstream, which is why the
validation pyramid of Section~\ref{sec:validationmap} tests the stages independently.

\paragraph{Dimensionalisation.} The entire eigenvalue pipeline is nondimensional. Physical units
enter only at the very end, through a \texttt{DimensionalEdgeState} $(U_e,\nu_e,T_e)$ and the
converters in \texttt{scales.py}: $F\to f$ [kHz], $R\to x$ [mm], $\sigma_L\to\sigma$ [1/m]. The
conversions apply the same $L^*$ definition the solver used, so no new physics and no new
approximation enters at this stage.

\paragraph{Geometry: the sharp cone.} For a sharp cone at zero incidence, the Mangler
transformation maps the cone boundary layer at surface distance $s$ onto the flat-plate problem
at $x_{\mathrm{eq}}=s/3$, independent of half-angle \cite{mack1984}; \texttt{pymack/cone.py} is
pure station bookkeeping ($R_{\mathrm{eq}}=R_s/\sqrt3$, $N_{\mathrm{cone}}=3N_{\mathrm{plate}}$
over the same $R_{\mathrm{eq}}$ window) and the solvers themselves are untouched. The cone path
therefore reuses stages 2--6 verbatim, changing only the interpretation of the abscissa.

%=====================================================================
\section{The canonical Mach 6 case, stage by stage}\label{sec:mach6}
%=====================================================================
The production case behind the README figures
(\texttt{scripts/run\_mach6\_spatial\_neutral\_case.py}; conditions $\Mae=6$, air, Sutherland,
$T_e=300$ K, $T_w/T_e=5.88$, isothermal wall) exercises the full map:
\begin{enumerate}[nosep,leftmargin=1.8em]
  \item \textbf{Stage 2.} One mean-flow solve at the stated wall and transport conditions.
  \item \textbf{Stage 3, formulation.} Spatial: the deliverable is a neutral curve in the
        $(R,F)$ plane and $N$-factors, both fixed-frequency products.
  \item \textbf{Stage 3, engine.} The dense spatial backend with continuation
        (\texttt{pymack\_dense} + \texttt{pymack\_continuation}), cross-checked against the main
        family by the consistency gate.
  \item \textbf{Stages 4--5.} One fixed-frequency sweep per $F\in[6.0\times10^{-5},
        2.30\times10^{-4}]$ over $R\in[200,3300]$, second-mode branch selected by the
        phase-speed window $0.90\le c\le0.97$; lower and upper neutral crossings extracted per
        frequency (policy: single sweep, no stitching, no smoothing, recorded in the run
        manifest).
  \item \textbf{Stage 6.} $N(R)$ and $e^N$ per frequency by Eq.~\eqref{eq:nfactor}; dimensional
        overlays via \texttt{DimensionalEdgeState}. Expected anchor: at $F=6.0\times10^{-5}$,
        $R_{\mathrm{lower}}\approx1543$, $R_{\mathrm{upper}}\approx3041$
        (\texttt{docs/MACH6\_SPATIAL\_NEUTRAL\_WORKFLOW.md}).
\end{enumerate}
Run \texttt{--geometry cone --cone-half-angle-deg 7} and the identical chain serves the cone
through the Mangler bookkeeping of Section~\ref{sec:deliverables}.

%=====================================================================
\section{Where validation attaches to the map}\label{sec:validationmap}
%=====================================================================
The layered validation strategy (\texttt{docs/VALIDATION\_STRATEGY.md}; \emph{Validation})
verifies the stages independently, so a failure localises to a stage instead of surfacing as an
unexplained error in a final product. Table~\ref{tab:valmap} lists, for each layer, the solver
route that produced the compared quantity, so every verdict can be traced back to a specific
engine of Section~\ref{sec:engines}:

\begin{table}[H]\centering\footnotesize
\renewcommand{\arraystretch}{1.35}
\begin{tabular}{@{}>{\raggedright\arraybackslash}p{0.7cm}>{\raggedright\arraybackslash}p{5.9cm}
>{\raggedright\arraybackslash}p{4.4cm}>{\raggedright\arraybackslash}p{3.2cm}@{}}
\toprule
\textbf{Layer} & \textbf{Stage verified: solver route (what it computed)} & \textbf{Benchmark} & \textbf{Status} \\
\midrule
0 & spectral machinery: $D_1,D_2$ and the domain map applied to analytic functions; no
    eigenvalue solver & analytic derivatives & machine precision \\
1 & stage 2: mean-flow similarity BVP (\texttt{baseflow.py}; \texttt{solve\_bvp}
    \cite{kierzenka2001}), thicknesses by quadrature & Mack Table 11.1 thicknesses
    \cite{mack1984} & $<0.5\%$ \\
2 & stages 3--4: Orr--Sommerfeld temporal GEVP by QZ (\texttt{solve\_temporal\_os},
    \texttt{solve\_poiseuille\_temporal}); least-stable eigenvalue & Orszag Poiseuille
    eigenvalue \cite{orszag1971} & 5+ significant figures \\
3 & stages 3--4: exact oblique first-order shooting, $6\times6$ and $8\times8$ systems
    (engine 3; \texttt{evaluate\_table\_10\_1\_exact\_shooting}); temporal growth rates
    $\omega_i$ & Mack Table 10.1 \cite{mack1984} & 0.07--0.91\% \\
4a & stage 3, spatial: five routes at one point, dense companion QEP
    (\texttt{dense.solve\_mack\_branch}) vs shift--invert QEP (\texttt{solve\_spatial}) vs
    Newton (\texttt{solve\_spatial\_newton}) vs Gaster-seeded pipeline
    (\texttt{solve\_spatial\_from\_temporal}) vs Muller (\texttt{solve\_spatial\_muller});
    complex $\alpha$ & internal cross-method redundancy & independent operators
    $\le2\times10^{-4}$ \\
4b & stage 3, spatial: shift--invert QEP (\texttt{solve\_spatial}); complex $\alpha$ at
    Malik's printed conditions & Malik (1990) Table IX Case 6 \cite{malik1990} & $\alpha$ to
    ${\sim}5\times10^{-6}$ \\
5 & stages 2--6 end to end: production fixed-frequency spatial sweep
    (\texttt{compute\_spatial\_fixed\_frequency\_curves.py}), neutral branches
    dimensionalised via \texttt{scales.py}; committed artifact with run manifest &
    independent-code Mach 5.35 N$_2$ neutral curve & upper branch MAE 3.2 mm \\
6 & overlays: \"Ozgen lobes from the temperature-form temporal GEVP
    (\texttt{solve\_temporal\_2d}); Mack Fig.~10.3 from the exact-shooting scan (engine 3) &
    \"Ozgen Fig.~3, Mack Fig.~10.3 \cite{ozgen2008,mack1984} & demonstrations, not gates \\
\bottomrule
\end{tabular}
\caption{The validation pyramid mapped onto the pipeline stages, with the solver route behind
each verdict. Tables and independent codes gate; figures illustrate.}
\label{tab:valmap}
\end{table}

\paragraph{The verification audit, case by case.} Beyond the CI gates, the standalone audit
(\texttt{verification/}\allowbreak\texttt{SUCCESS\_MATRIX.md}) runs published benchmark conditions and records a
verdict per case. Table~\ref{tab:auditmap} states which solver produced each compared quantity;
the pattern is the division of labour of Sections~\ref{sec:formulations}--\ref{sec:engines} in
action: temporal GEVP sweeps for growth-rate figures, the spatial QEP family for neutral
branches, eigenvalue anchors, and $N$-factors, and the exact shooting engine wherever the
reduced EVP's oblique first-mode branch selection is not trusted.

\begin{table}[H]\centering\footnotesize
\renewcommand{\arraystretch}{1.35}
\begin{tabular}{@{}>{\raggedright\arraybackslash}p{3.4cm}>{\raggedright\arraybackslash}p{5.6cm}
>{\raggedright\arraybackslash}p{3.0cm}>{\raggedright\arraybackslash}p{2.2cm}@{}}
\toprule
\textbf{Case} & \textbf{Solver route} & \textbf{Quantity compared} & \textbf{Verdict} \\
\midrule
Mack Fig.~10.6, $M$=4.5--10 \cite{mack1984} & temporal GEVP sweep over $\alpha$
  (\texttt{solve\_temporal\_compressible}) & second-mode max $\omega_i$ vs $M$ & 1--7\% \\
Ma \& Zhong (2003), $M$=4.5 \cite{mazhong2003} & spatial QEP fixed-frequency sweeps
  (\texttt{solve\_spatial}), $\sigma=0$ crossings & second-mode neutral branches I/II &
  ${\sim}3\%$ \\
Balakumar \& Malik (1992) \cite{balakumar1992} & full companion-QEP spectrum
  (\texttt{solve\_spatial\_full\_spectrum}) & discrete eigenvalue + discrete/continuous branch
  topology & $\alpha_r$ exact; $\alpha_i$ ${\sim}10\%$ \\
Egorov et al.\ (2006), $M$=6 & spatial QEP growth curves (\texttt{solve\_spatial}) &
  most-amplified frequency vs DNS-LST & ${\sim}7\%$ \\
Sivasubramanian \& Fasel (2015), $M$=6 cone & spatial QEP (\texttt{solve\_spatial}) + Mangler
  bookkeeping (\texttt{cone.py}) + trapezoid $N$; Taylor--Maccoll edge state & sharp-cone
  $N$-factor & $N\approx7.1$ vs ${\sim}7$--$8$ \\
Malik (1990) anchors \cite{malik1990} & Cases 3--5: temporal GEVP
  (\texttt{solve\_temporal\_compressible}, \texttt{\_3d}); Case 6, Table X: spatial QEP
  (\texttt{solve\_spatial}); Fig.~4: temporal eigenfunction & eigenvalue anchors +
  second-mode eigenfunction & anchors within literature spread \\
Mack Fig.~10.1, $M$=2.2 \cite{mack1984} & temporal GEVP neutral tracing
  (\texttt{solve\_temporal\_compressible}) & first-mode neutral loop & documented weak spot \\
Mack Fig.~10.3 \cite{mack1984} & exact 3D shooting continuation from an anchor root
  (\texttt{temporal\_growth\_scan\_3d\_shooting\_from\_anchor}) & oblique first-mode
  $\omega_i(R)$ & open ($M{\approx}2.2$) \\
Mack Fig.~10.4 \cite{mack1984} & reduced 3D temporal GEVP
  (\texttt{solve\_temporal\_compressible\_3d}) & high-$M$ oblique first mode & under-amplified \\
\"Ozgen Fig.~3, $M$=2--8 \cite{ozgen2008} & temperature-form temporal GEVP
  (\texttt{solve\_temporal\_2d}) + discrete-mode acceptance + branch continuation & first/second
  mode growth lobes & mixed (second mode agrees; first mode under-amplified) \\
\bottomrule
\end{tabular}
\caption{The verification-audit cases with the solver route behind each verdict. Verdicts
paraphrase \texttt{verification/SUCCESS\_MATRIX.md}, which remains the single source of truth;
each case folder's \texttt{verdict.json} records metrics and provenance.}
\label{tab:auditmap}
\end{table}

A confidence note: the second (Mack) mode, the design target, is validated across
Mach 4.5--10 against Mack, Ma \& Zhong, Balakumar \& Malik, Egorov et al., and a sharp-cone
$N$-factor comparison \cite{mack1984,mazhong2003,balakumar1992}; the first mode at low-to-moderate
Mach is systematically under-amplified in the reduced EVP path and is the documented weak spot,
partially recovered by the exact shooting engine (Section~\ref{sec:engines}, item 3).

\begin{thebibliography}{99}
\bibitem{mack1984} L.~M. Mack, ``Boundary-layer linear stability theory,'' AGARD Report No.~709,
  \emph{Special Course on Stability and Transition of Laminar Flow}, 1984.
\bibitem{malik1990} M.~R. Malik, ``Numerical methods for hypersonic boundary layer stability,''
  \emph{Journal of Computational Physics}, Vol.~86, No.~2, 1990, pp.~376--413.
\bibitem{ozgen2008} S.~\"Ozgen and S.~A. K{\i}rcal{\i}, ``Linear stability analysis in
  compressible, flat-plate boundary layers,'' \emph{Theoretical and Computational Fluid
  Dynamics}, Vol.~22, 2008, pp.~1--20.
\bibitem{orszag1971} S.~A. Orszag, ``Accurate solution of the Orr--Sommerfeld stability
  equation,'' \emph{Journal of Fluid Mechanics}, Vol.~50, 1971, pp.~689--703.
\bibitem{schmid2001} P.~J. Schmid and D.~S. Henningson, \emph{Stability and Transition in Shear
  Flows}, Springer, 2001.
\bibitem{molerstewart1973} C.~B. Moler and G.~W. Stewart, ``An algorithm for generalized matrix
  eigenvalue problems,'' \emph{SIAM Journal on Numerical Analysis}, Vol.~10, 1973, pp.~241--256.
\bibitem{tisseur2001} F.~Tisseur and K.~Meerbergen, ``The quadratic eigenvalue problem,''
  \emph{SIAM Review}, Vol.~43, No.~2, 2001, pp.~235--286.
\bibitem{gaster1962} M.~Gaster, ``A note on the relation between temporally-increasing and
  spatially-increasing disturbances in hydrodynamic stability,'' \emph{Journal of Fluid
  Mechanics}, Vol.~14, 1962, pp.~222--224.
\bibitem{trefethen2000} L.~N. Trefethen, \emph{Spectral Methods in MATLAB}, SIAM, 2000.
\bibitem{conte1966} S.~D. Conte, ``The numerical solution of linear boundary value problems,''
  \emph{SIAM Review}, Vol.~8, No.~3, 1966, pp.~309--321.
\bibitem{neldermead1965} J.~A. Nelder and R.~Mead, ``A simplex method for function
  minimization,'' \emph{The Computer Journal}, Vol.~7, No.~4, 1965, pp.~308--313.
\bibitem{smith1956} A.~M.~O. Smith and N.~Gamberoni, ``Transition, pressure gradient and
  stability theory,'' Douglas Aircraft Co.\ Report ES-26388, 1956.
\bibitem{vaningen1956} J.~L. van Ingen, ``A suggested semi-empirical method for the calculation
  of the boundary layer transition region,'' Delft University of Technology, Report VTH-74,
  1956.
\bibitem{kierzenka2001} J.~Kierzenka and L.~F. Shampine, ``A BVP solver based on residual
  control and the MATLAB PSE,'' \emph{ACM Transactions on Mathematical Software}, Vol.~27,
  No.~3, 2001, pp.~299--316.
\bibitem{mazhong2003} Y.~Ma and X.~Zhong, ``Receptivity of a supersonic boundary layer over a
  flat plate. Part 1. Wave structures and interactions,'' \emph{Journal of Fluid Mechanics},
  Vol.~488, 2003, pp.~31--78.
\bibitem{balakumar1992} P.~Balakumar and M.~R. Malik, ``Discrete modes and continuous spectra in
  supersonic boundary layers,'' \emph{Journal of Fluid Mechanics}, Vol.~239, 1992, pp.~631--656.
\end{thebibliography}

\end{document}
